\ExerciseSection

\begin{enumerate}
\def\labelenumi{\arabic{enumi})}
\item
  设 X 是一个豪斯多夫一致空间，\((f_{t})\) 是 X 上定义 X 的一致结构的一族伪度量。设 \(X_{t}\) 是这样的一致空间相联系的度量空间：该一致空间由 X 赋以单个伪度量 \(f_{t}\) 所定义的结构而得（第 1 号）。证明 X 同构于乘积一致空间 \(\prod_{i} X_{i}\) 的一个子空间。
\item
  在一个连通的度量空间 X 中，若度量在 \(X \times X\) 上不是有界的，证明球面不可能是空的。
\end{enumerate}

¶ 3) a) 设 X 是一个紧度量空间。若每个开球在 X 中的闭包都是同中心同半径的闭球，则 X 中每个球都是连通集{[}若 x 与 y 是 X 的两点，S 是以 x 为中心、\(d(x, y)\) 为半径的闭球，证明对每个 \(\varepsilon > 0\) ，点所成的集 \(A_{y,\varepsilon}\) ——S 中可用一条含于 S 的 \(V_{\varepsilon}\) 链与 y 相连的那些点——（第 II 章，§ 4，第 4 号）含有满足 \(d(x, z) < d(x, y)\) 的点 z {]}。用反证法并利用魏尔斯特拉斯定理（第 IV 章，§ 6，第 1 号，定理 1），推出 \(x \in A_{y,\varepsilon}\) 对所有 \(\varepsilon > 0\) 成立。

\begin{enumerate}
\def\labelenumi{\alph{enumi})}
\setcounter{enumi}{1}
\tightlist
\item
  在空间 \(\mathbf{Y} = \mathbf{R}^2\) （赋以度量
\end{enumerate}

\[
d (x, y) = \max \left(\left| x _ {1} - y _ {1} \right|, \left| x _ {2} - y _ {2} \right|\right),
\]

）中，设 X 是由所有满足下列条件的点 \((x_{1}, x_{2})\) 组成的紧子空间：或者 \(x_{1} = 0\) 且 \(0 \leqslant x_{2} \leqslant 1\) ，或者 \(0 \leqslant x_{1} \leqslant 1\) 且 \(x_{2} = 0\) 。证明 X 中每个球都连通，但一个开球的闭包未必是同中心同半径的闭球。

\begin{enumerate}
\def\labelenumi{\arabic{enumi})}
\setcounter{enumi}{3}
\tightlist
\item
  度量空间 \(\mathbf{X}\) 称为超度量空间，如果它的度量 \(d\) 满足不等式
\end{enumerate}

\[
d (x, y) \leqslant \sup \left(d (x, z), d (y, z)\right)
\]

（对所有 \(x, y, z\) ∈ \(\mathbf{X}\) ；该不等式蕴含三角不等式）（参看 § 6，习题 2）。

\begin{enumerate}
\def\labelenumi{\alph{enumi})}
\item
  若 \(d(x, z) \neq d(y, z)\) ，证明 \(d(x, y) = \sup (d(x, z), d(y, z))\) 。
\item
  设 \(\mathrm{V}_{r}(x)\) 是以 x 为中心、r 为半径的开球。证明 \(\mathrm{V}_{r}(x)\) 在 X 中既开又闭（因而 X 是全不连通的），且对每个 \(y \in \mathrm{V}_{r}(x)\) 有 \(\mathrm{V}_{r}(y) = \mathrm{V}_{r}(x)\) 。
\item
  证明以 x 为中心、r 为半径的闭球 \(\mathrm{W}_{r}(x)\) 在 X 中既开又闭，且对每个 \(y \in \mathrm{W}_{r}(x)\) 有
\end{enumerate}

\[
\mathrm{W} _ {r} (y) = \mathrm{W} _ {r} (x).
\]

含于 \(W_{r}(x)\) 中的、半径为 r 的各个不同开球构成 \(W_{r}(x)\) 的一个分划，且这些球中任何两个之间的距离都等于 r 。d) 若 X 的两个球（开的或闭的）相交，则其中一个含于另一个之中。

\begin{enumerate}
\def\labelenumi{\alph{enumi})}
\setcounter{enumi}{4}
\item
  X 的点序列 \((x_{n})\) 是柯西序列，当且仅当 \(d(x_{n}, x_{n+1})\) 当 \(n\) 趋于无穷时趋于 0 。
\item
  若 X 是紧的，证明对每个 \(x_{0} \in X\) ，\(d(x_{0}, x)\) 在 X 中所取值的集合是 \([0, +\infty]\) 的一个（有限或无限的）可数子集，其所有点除 0 可能例外外都是孤立的{[}对 \(d(x_{0}, x)\) 所取的每个值 r，考虑 \(d(x_{0}, x)\) 在满足 \(d(x_{0}, x) < r\) 的点集上的上确界，以及它在满足 \(d(x_{0}, x) > r\) 的点集上的下确界{]}。
\end{enumerate}

\begin{enumerate}
\def\labelenumi{\arabic{enumi})}
\setcounter{enumi}{4}
\tightlist
\item
  设 X 是一个度量空间，d 是它的度量。对 X 的每一对点 \((x, y)\) ，令 \(d_{0}(x, y)\) 表示满足如下条件的实数 \(\alpha > 0\) 组成的集合的下确界：x 与 y 可以用一条 \(V_{\alpha}\) 链连接起来（第 II 章，§ 4，第 4 号）。证明 \(d_{0}\) 是 X 上的一个伪度量，且 X 上由伪度量 \(d_{0}\) 所定义的一致空间（第 1 号）相联系的度量空间是一个超度量空间（习题 4）。
\end{enumerate}

§ 6) 设 X 是一个度量空间。若 A 与 B 是 X 的两个非空子集，令

\[
\rho (\mathrm{A}, \mathrm{B}) = \sup _ {x \in \mathbf {A}} d (x, \mathrm{B}) \quad \text { and } \quad \sigma (\mathrm{A}, \mathrm{B}) = \sup \left(\rho (\mathrm{A}, \mathrm{B}), \rho (\mathrm{B}, \mathrm{A})\right);
\]

又令 \(\sigma(\emptyset, \emptyset) = 0\) ，\(\sigma(\emptyset, A) = \sigma(A, \emptyset) = +\infty\) （\(A\) 为 \(X\) 中任意非空子集）。证明 \(\sigma\) 是集合 \(\mathfrak{P}(X)\) （即 \(X\) 的全体子集组成的集合）上的一个伪度量，且由 \(\sigma\) 定义的一致结构与按第 II 章，§ 1，习题 5 的程序由 \(X\) 的一致结构构造出的一致结构相重合。

假设 X 是有界的。此时 \(\sigma\) 是集合 \(\mathfrak{F}(\mathbf{X})\) （X 的非空闭子集全体）上的一个度量。此外若 X 还是完备的，证明 \(\mathfrak{F}(\mathbf{X})\) 是一个完备度量空间{[}设 \(\Phi\) 是 \(\mathfrak{F}(\mathbf{X})\) 上的一个柯西滤子；对每个集 \(\mathcal{X} \in \Phi\) ，令 \(S(\mathcal{X})\) 表示属于 \(\mathcal{X}\) 的那些 X 的子集 A 的并集；证明诸集 \(S(\mathcal{X})\) 当 \(\mathcal{X}\) 跑遍 \(\Phi\) 时构成 X 上的一个滤子基，该滤子基具有一个非空的聚点集 C，且 \(\Phi\) 收敛于 C{]}。

\begin{enumerate}
\def\labelenumi{\arabic{enumi})}
\setcounter{enumi}{6}
\item
  设 X 是一个无限离散空间。证明 X 上的有限分划的一致结构（第 II 章，§ 2，第 2 号）——它与 X 的拓扑相容——不是可度量化的（否则存在 X 的有限分划的序列 \((\mathfrak{F}_n)\) ，使得 X 的每个有限分划都由一些集组成，其中每个集是属于诸分划 \(\mathfrak{F}_n\) 之一的集的并；由此将推出 X 的有限分划组成的集合是可数的）。
\item
  设 \((\mathbf{X}_{t})\) 是一个不可数的豪斯多夫拓扑空间族，其中每个空间至少有两个不同的点。证明在乘积空间 \(\prod_{i} X_{i}\) 中，任何点都没有可数的基本邻域系。
\item
  设 X 是一个拓扑空间，它的每个点都有可数的基本邻域系。
\end{enumerate}

\begin{enumerate}
\def\labelenumi{\alph{enumi})}
\item
  只要 X 中每个收敛序列都只有一个极限，X 就是豪斯多夫的。
\item
  若 a 是 X 的点组成的无限序列 \((x_{n})\) 的一个聚点，则存在 \((x_{n})\) 的一个无限子序列收敛于 a。
\item
  设 A 是 X 的一个非空子集，\(x_{0}\) 是 \(\overline{A}\) 的一个点，f 是 A 到一个豪斯多夫拓扑空间 \(X'\) 中的映射。则点 \(a \in X'\) 是 f 在点 \(x_{0}\) 处（关于 A）的极限，只要对 A 中点的每个序列 \((x_{n})\) （它收敛于 \(x_{0}\) ），序列 \((f(x_{n}))\) 都收敛于 a。
\item
  沿用 c) 的记号，再设 \(X'\) 的每个点都有可数的基本邻域系。若 \(a \in X'\) 是 f 在 \(x_{0}\) 处关于 A 的一个聚点，则存在 A 的点组成的序列 \((x_{n})\) ，它收敛于 \(x_{0}\) ，且使序列 \((f(x_{n}))\) 收敛于 a。
\end{enumerate}

¶ 10) 设 X 是一个紧空间。若存在 X×X 上的一个连续实值函数 f，使得关系 \(f(x,y)=0\) 等价于 x=y，则 X 是可度量化的。（证明若 V 跑遍 R 中 0 的一个基本邻域系，则诸集 \(\overline{f}^{1}(V)\) 构成 X×X 中对角线 \(\Delta\) 的一个基本邻域系。）

¶ 11) 设 X 是一个紧度量空间，d 是它的度量。证明：若 f 是 X 到 X 中的映射，使得对 X 的每一对点 x、y 有 \(d(f(x), f(y)) \geqslant d(x, y)\) ，则 f 是 X 到其自身上的一个等距映射。{[}设 a、b 是 X 的任意两点，令 \(f^n = f^{n-1} \circ f\) ，\(a_n = f^n(a)\) ，\(b_n = f^n(b)\) ；证明对每个 \(\varepsilon > 0\) ，存在指标 k 使得 \(d(a, a_k) \leqslant \varepsilon\) 且 \(d(b, b_k) \leqslant \varepsilon\) （选取 \((a_n)\) 与 \((b_n)\) 的适当子序列）；由此推出 \(d(a_1, b_1) = d(a, b)\) ，且 \(f(\mathbf{X})\) 在 X 中稠密。{]}

¶ 12) 设 X 是一个紧可度量化空间。

\begin{enumerate}
\def\labelenumi{\alph{enumi})}
\item
  证明存在康托尔三分集 K（第 IV 章，§ 2，第 5 号）到 X 上的一个连续映射（参看第 IV 章，§ 8，习题 11）。
\item
  若此外 X 还是全不连通的且没有孤立点，则 X 同胚于 K。{[}按第 IV 章，§ 8，习题 12 的方式论证，利用如下事实：点 \(x \in \mathbf{X}\) 的每个邻域都包含 \(x\) 的一个既开又闭的邻域（第 II 章，§ 4，第 4 号，命题 6 的推论）{]}。
\end{enumerate}

¶ 13) a) 在实数集 R 上，考虑如下定义的拓扑 C：对每个 y \textgreater{} 0，\(\mathrm{U}_{y}(x)\) 表示诸区间的并 \([x, x + y[ \text{ 和 } ] - x - y, -x[; C \text{ 是满足下述条件的拓扑： } \mathrm{U}_{y}(x) \text{ 构成邻域基本系，其中心为 } x \text{ 当 } y \text{ 遍历实数集 } > 0. \text{ 令 X 为在区间 } [-1, +1] \text{ 上赋予由 C 诱导的拓扑所得到的空间。证明 X 是紧的。（考虑 X 上的一个滤子 F；若 } x \in X \text{ 是 F 在实直线拓扑下的聚点，证明 } x \text{ 或 } -x \text{ 是 F 在 C 下的聚点.)}\)

\begin{enumerate}
\def\labelenumi{\alph{enumi})}
\setcounter{enumi}{1}
\item
  X 的每个点都有可数的基本邻域系，且 X 有可数稠密子集，但 X 不是可度量化的（证明它的拓扑没有可数基）。
\item
  设 A 是 X 的一个开子集。证明 A 是如下诸集之并：开区间组成的一个可数族 \((I_{n})\) （这些区间含于 {[}0, 1{]} ）、诸区间 \(-I_{n}\) 、诸区间 \(I_{n}\) 与 \(-I_{n}\) 的左端点组成的集合的一个子集，以及可能还有点 +1。由此证明 A 是 X 的闭子集组成的一个可数族之并。
\item
  设 Y 是集合 \(I \times \{1, 2\}\) ，其中 I 是 R 的区间 {[}0, 1{]} ；令 \(X_{i}\) 表示 \(I \times \{i\} (i = 1, 2)\) ，令 f 表示双射 \((x, 1) \to (x, 2)\) （它是 \(X_{1}\) 到 \(X_{2}\) 上的映射）。对每个 \(x \in I\) ，令 \(\mathfrak{B}((x, 2))\) 表示仅由 \(\{(x, 2)\}\) 组成的子集所成的集合，令 \(\mathfrak{B}((x, 1))\) 表示形如 \(V_{1} \cup (f(V_{1}) - \{x, 2\})\) 的子集所成的集合，其中 \(V_{1} = V \times \{1\}\) 而 V 跑遍 I 中 x 的一个基本邻域系。证明对每个 \(y \in Y\) ，\(\mathfrak{B}(y)\) 是 Y 上某个拓扑的、y 的基本邻域系。赋以此拓扑后，Y 是紧的，且 Y 的每个点都有可数的基本邻域系，但 Y 没有可数稠密子集。Y 的每个含于 \(X_{2}\) 中的闭子集都是有限的；由此证明 Y 的紧子集 \(X_{1}\) 没有可数的基本邻域系。若 R 是 Y 上的等价关系，其等价类是集合 \(X_{1}\) 以及 \(Y - X_{1}\) 的诸点，则 R 是闭的，且模 R 的每个等价类都是紧的，但 Y/R 中有一个点没有可数的基本邻域系{[}参看 § 4，习题 24 a){]}。
\end{enumerate}

\begin{enumerate}
\def\labelenumi{\arabic{enumi})}
\setcounter{enumi}{13}
\tightlist
\item
  \begin{enumerate}
  \def\labelenumii{\alph{enumii})}
  \tightlist
  \item
    设 X 是一个可达的拓扑空间（第 I 章，§ 8，习题 1）。证明下列性质等价：
  \end{enumerate}
\end{enumerate}

α) X 的每个点序列都有聚点。

\(\beta)\) \(\mathbf{X}\) 的任何无限离散子空间都不是闭的。

\(\gamma)\) \(\mathbf{X}\) 的每个可数开覆盖都包含 \(\mathbf{X}\) 的一个有限开覆盖。

δ) 对 X 的任何无限开覆盖 R，存在 X 的一个开覆盖 S ⊂ R，它异于 R。

拓扑空间 X 称为可数紧的，如果它是豪斯多夫的且具有上述性质。

\begin{enumerate}
\def\labelenumi{\alph{enumi})}
\setcounter{enumi}{1}
\item
  可数紧空间中的点序列是收敛的，当且仅当它有唯一的聚点。
\item
  可数紧空间的每个闭子空间都是可数紧的。反之，若 X 是豪斯多夫的且 X 的每个点都有可数的基本邻域系，则 X 的每个可数紧子空间在 X 中是闭的。
\item
  设 \(f\) 是可数紧空间 \(\mathbf{X}\) 到豪斯多夫空间 \(\mathbf{X}'\) 中的一个连续映射。则 \(f(\mathbf{X})\) 是 \(\mathbf{X}'\) 的一个可数紧子空间。
\item
  设 X 是一个可数紧空间，它的每个点都有可数的基本邻域系。则 X 的每个点序列都有收敛子序列。
\end{enumerate}

\(f\) ) 设 \((\mathbf{X}_n)\) 是拓扑空间组成的一个可数序列，其中每个空间的每个点都有可数的基本邻域系。

此时乘积空间 \(\prod_{n=1}^{\infty} \mathbf{X}_n\) 是可数紧的，当且仅当诸空间 \(\mathbf{X}_n\) 中的每一个都是可数紧的。{[}利用 e) 及第 I 章，§ 6，习题 16。{]}

\begin{enumerate}
\def\labelenumi{\alph{enumi})}
\setcounter{enumi}{6}
\item
  每个点都有可数基本邻域系的可数紧空间是正则的。
\item
  若一个可数紧空间具有可数的开集基，则它是紧的，从而是可度量化的。
\item
  可数紧空间上的每个下半连续实值函数都达到它的下确界。特别地，每个可数紧空间都是伪紧的（§ 1，习题 21）{[}参看 § 4，习题 26 b){]}。
\item
  证明性质 \(\delta\) （见 \(a\) ）蕴含下列性质：
\end{enumerate}

ζ) X 的每个逐点有限（§ 4，第 3 号）开覆盖都包含 X 的一个有限开覆盖。（用反证法论证。）反之，满足 ζ) 的正则空间是可数紧的{[}证明此时 ζ) 蕴含 β){]}。

¶ 15) 设 \(\mathbf{X} = [a, b[\) 是第 I 章，§ 9，习题 12 中定义的局部紧空间。

\begin{enumerate}
\def\labelenumi{\alph{enumi})}
\item
  X 的一个子集是相对紧的，当且仅当它是有界的。由此证明 X 是可数紧的，从而（命题 15）不是可度量化的（注意 X 的每个可数子集都是有界的）。
\item
  证明 X 的每个点都有一个可度量化邻域（利用命题 16）。
\item
  若 A 与 B 是 X 中两个非紧的闭集，则它们的交不是空的（构造一个递增序列，使其偶数编号的项是 A 的点，奇数编号的项是 B 的点）。由此推出，非紧闭集的每个邻域都是一个相对紧集的余集。若 A 是 X 的非孤立点组成的集合，证明 A 是闭的，且不是任何可数开集族的交。
\item
  令 \(\mathbf{X}'\) 表示赋有如下拓扑的区间 \([a, b]\) ：对每个 \(x \in \mathbf{X}\) ，诸区间 \(]y, x]\) （其中 \(y\) 跑遍由满足 \(< x\) 的元素组成的集合）构成 \(x\) 的一个基本邻域系；\(b\) 的一个基本邻域系由诸集 \(\mathbf{V}_x\) 组成，其中对每个 \(x \in \mathbf{X}\) ，\(\mathbf{V}_x\) 表示 \(\{b\}\) 与 \([x, b]\) 中具有前元的点（即在 \(\mathbf{X}\) 中孤立的那些点）组成的集合之并。证明 \(\mathbf{X}'\) 是可数紧的但不是正则的{[}参看习题 14 g){]}，且子空间 \(\mathbf{X}\) （它含于 \(\mathbf{X}'\) ）是可数紧的但不是闭的{[}参看习题 14 c){]}。
\end{enumerate}

\begin{enumerate}
\def\labelenumi{\arabic{enumi})}
\setcounter{enumi}{15}
\tightlist
\item
  设 X 与 Y 是两个可数紧空间。证明在下列两种情形的每一种中，乘积 \(X \times Y\) 都是可数紧的：(i) X、Y 之一是紧的；(ii) X、Y 之一使每个点都有可数的基本邻域系。{[}在情形 (i)，设 Y 是紧的，考虑一个递增的开集序列 \((\mathbf{G}_{n})\) ，其并为 \(X \times Y\) ；对每个 n，令 \(H_{n}\) 表示所有满足如下条件的 \(x \in X\) 组成的集合：\(\{x\} \times Y \subset G_{n}\) ；证明 \(H_{n}\) 是开的，且 \(X = \bigcup_{n} H_{n} \cdot\) {]}
\end{enumerate}

¶ 17) 设 βN 是离散空间 N 的斯通–切赫紧化（§ 1，习题 7）。

\begin{enumerate}
\def\labelenumi{\alph{enumi})}
\item
  证明存在 \(\beta N\) 的两个可数紧子空间 A、B，使得 \(A \cap B = N\) 且 \(A \cup B = \beta N\) 。{[}令 \(\aleph_{\alpha} = 2^{2^{\operatorname{Card}(N)}}\) ；根据 § 1，习题 12 b)，存在一个双射 \(\xi \to S_{\xi}\) ，它把序数 \(\omega_{\alpha}\) （《集合论》R, 第 III 章，§ 6，习题 10）映到 X 的可数无限子集组成的集合上。用超限归纳法定义两个单射映射 \(\xi \to x_{\xi}, \xi \to y_{\xi}\) ，它们把 \(\omega_{\alpha}\) 映入 \(\beta N - N\) ，使得 \(x_{\xi} \in \overline{S}_{\xi}, y_{\xi} \in \overline{S}_{\xi}\) 对每个 \(\xi\) 成立，且使得若 P（相应地 Q）是诸 \(x_{\xi}\) （相应地诸 \(y_{\xi}\) ）组成的集合，则有 \(P \cap Q = \varnothing\) ，\(P \cup Q = \beta N - N\) ；为此，利用 § 1，习题 12 b) 与 c)。然后证明 A = P ∪ N 与 B = Q ∪ N 满足问题的条件。{]}
\item
  证明乘积空间 \(A \times B\) 不是可数紧的（注意 \(A \times B\) 与 \(\beta N \times \beta N\) 的对角线的交是 \(A \times B\) 的一个无限离散闭子空间）。
\end{enumerate}

¶ 18) 设 X 是一个度量空间，使得对每个 \(x \in X\) ，存在一个以 \(x\) 为中心的开球，它作为 X 的子空间考虑时具有可数基。令 \(r_x\) 表示具有该性质的、以 \(x\) 为中心的球的半径所组成的集合的上确界。

\begin{enumerate}
\def\labelenumi{\alph{enumi})}
\item
  利用佐恩引理证明：存在 X 中两两不相交的开球组成的一个极大族 \((\mathbf{B}_{\alpha})\) ，使得若 \(x_{\alpha}\) 是 \(B_{\alpha}\) 的中心，则 \(B_{\alpha}\) 的半径 \(<r_{x_{\alpha}}\) 。证明诸 \(B_{\alpha}\) 的并在 X 中稠密，从而存在 X 的一个稠密子集 M，使得以任意点 \(x \in X\) 为中心、半径 \(<r_{x}\) 的每个开球只含 M 的可数无穷多个点（注意这样一个球只能与可数无穷多个两两不相交的开集相交，并利用命题 12）。
\item
  对每个点 \(x \in \mathbf{M}\) ，令 \(\mathbf{S}_x\) 表示以 \(x\) 为中心、\(\frac{1}{3} r_x\) 为半径的开球。证明诸 \(\mathbf{S}_x\) 覆盖 \(\mathbf{X}\) ，且 \(\mathbf{X}\) 的任意点只属于可数无穷多个 \(\mathbf{S}_x\) {[}注意若 \(y \in \mathbf{S}_x\) ，则有 \(d(x, y) \leqslant \frac{1}{2} r_y\) {]}。
\item
  设 R 是 X 的点 \(x, y\) 之间的如下等价关系：存在 M 的点组成的序列 \((z_i)_{1 \leqslant i \leqslant n}\) ，使得 \(x \in S_{z_i}\) ，\(y \in S_{z_n}\) ，且 \(S_{z_i}\) 与 \(S_{z_{i+1}}\) 相交（对 \(1 \leqslant i \leqslant n - 1\) ）。证明（模 R 的）等价类是 X 的一些子空间，它们在 X 中既开又闭，且具有可数基；换言之，X 是具有可数基的一些度量空间的拓扑和（第 I 章，§ 2，第 4 号）。
\end{enumerate}

¶ 19) 在度量空间中，每个相对紧集都是有界的。若 X 是一个拓扑空间，证明 X 的下列条件等价：α) 存在 X 上的一个度量，它与 X 的拓扑相容，使得 X 的每个（关于该度量的）有界子集都是相对紧的；β) X 是局部紧的且具有可数基。{[}为证 α) → β)，证明若每个有界集都相对紧，则 X 是局部紧的且是 σ-紧的。为证 β) → α)，注意由命题 12 的推论，X 是可度量化的；若 d 是 X 上与 X 的拓扑相容的一个度量，且 f 是 § 1，习题 15 d) 中定义的函数，则取 X 上由两个伪度量 d(x, y) 与 \textbar f(x) --- f(y)\textbar 所定义的一致结构。{]}

¶ 20) a) 给出一个例子：在具有可数基的可度量化局部紧空间 X 上的一个闭等价关系 R，使得 X/R 是仿紧的，但 X/R 中存在一个点没有可数的基本邻域系{[}参看第 I 章，§ 10，习题 17 及第 IX 章，§ 4，习题 24 a){]}。

\begin{enumerate}
\def\labelenumi{\alph{enumi})}
\setcounter{enumi}{1}
\tightlist
\item
  设 X 是一个可度量化空间，R 是 X 上的一个闭等价关系。证明：若 X/R 的每个点都有可数的基本邻域系，则每个（模 R 的）等价类在 X 中具有紧的边界。（证明：若结论不真，则存在 X 中的一个序列，其点两两不同、没有聚点，而它在 X/R 中的像是一个序列，收敛到与该序列所有点都不同的一个点。）对每个 \(z \in X/R\) ，令 \(C_{z}\) 为 z 在 X 中的逆像，且令 \(F_{z}\) 为 \(C_{z}\) 的边界（若该边界非空）；若 \(C_{z}\) 既开又闭，令 \(F_{z}\) 为由 \(C_{z}\) 的单独一个点组成的任一子集。令 Y 为当 z 跑遍 X/R 时诸集 \(F_{z}\) 的并。证明 \(Y/R_{y}\) 同胚于 X/R。
\end{enumerate}

¶ 21) a) 设 X 是一个可度量化空间，d 是与 X 的拓扑相容的一个有界度量，σ 是 d 在非空闭子集组成的集合 ℱ(X) 上对应的度量（习题 6）；R 是 X 上的一个既开又闭的等价关系。证明 σ 在 X/R 上的限制是与商拓扑相容的一个度量。{[}利用习题 20 b) 证明模 R 的每个等价类都是开的且紧的。然后证明：若 X/R 中的序列 (zn) 趋于一个点 a，且 φ: X → X/R 是典范映射，则有 σ(φ̄¹(a), φ̄¹(zₙ)) → 0；利用 R 是开的这一事实，用反证法论证。{]}

\begin{enumerate}
\def\labelenumi{\alph{enumi})}
\setcounter{enumi}{1}
\tightlist
\item
  在 R 的紧区间 I = {[}0, 1{]} 上，设 R 是这样的等价关系：其等价类为康托尔集 K（第 VI 章，§ 2，第 5 号）中除 K 的邻接区间的端点以外的点，以及 K 的邻接区间的闭包。证明商空间 I/R 同胚于 I{[}第 IV 章，§ 8，习题 16 b){]}，但度量 \(\sigma\) 与 I/R 的拓扑不相容。
\end{enumerate}

\begin{enumerate}
\def\labelenumi{\arabic{enumi})}
\setcounter{enumi}{21}
\tightlist
\item
  设 X 是具有可数基 \((\mathbf{U}_{n})\) 的一个豪斯多夫拓扑空间，R 是 X 上的一个豪斯多夫等价关系，使得 X/R 的每个点都有可数的基本邻域系。证明 X/R 的拓扑具有可数基。{[}设 \(\varphi: X \to X/R\) 是典范映射，证明诸集 \(\varphi(\mathbf{U}_{n})\) 的有限并的内部构成 X/R 的拓扑的一个基：若 V 是点 \(z \in X/R\) 的一个邻域，且 \((\mathbf{W}_{k})\) 是属于基 \((\mathbf{U}_{n})\) 、含于 \(\overline{\varphi}^{1}(\mathbf{V})\) 且覆盖 \(\overline{\varphi}^{1}(z)\) 的集组成的一个序列，证明存在有限多个指标 k，使诸 \(\varphi(\mathbf{W}_{k})\) 的并是 z 的一个邻域。用反证法：若此命题不真，则可构造 X/R 的两两不同的点组成的序列 \((y_{n})\) ，它趋于 z，且使 \(y_{n}\) 不属于诸 \(\varphi(\mathbf{W}_{k})\) （\(k \leqslant n\) ）的并；然后证明诸 \(\overline{\varphi}^{1}(y_{n})\) 的并在 X 中是闭的。{]}
\end{enumerate}

证明若假设 R 是闭的且模 R 的每个等价类都是紧的，则同样的结论成立（方法类似）。

¶ 23) 拓扑空间称为次可度量化的，如果它的拓扑比某个可度量化空间的拓扑更细。次可度量化空间是豪斯多夫的，但未必是正则的（第 I 章，§ 8，习题 20）。

\begin{enumerate}
\def\labelenumi{\alph{enumi})}
\item
  完全正则空间 X 是次可度量化的，当且仅当存在 X 上的一个一致结构，它与 X 的拓扑相容，且由一个伪度量族 \(\Phi\) 定义，该族至少包含一个度量 d（参看 § 1，习题 3）。设 \(\hat{X}\) 是 X 关于这个一致结构的完备化；度量 d 可扩张为一个伪度量 \(\overline{d}\) （在 \(\hat{X}\) 上）；证明对每个 \(x \in \hat{X}\) ，至多存在一个点 \(y \in X\) 使得 \(\overline{d}(x, y) = 0\) 。
\item
  证明若 X 是一个完全正则的次可度量化空间，则存在与 X 的拓扑相容的一个一致结构，X 关于它是完备的{[}利用 a) 及 § 1，习题 16 c){]}。由此证明若此外 X 还是伪紧的（§ 1，习题 21），则 X 是紧的。
\item
  证明局部紧、仿紧、次可度量化的空间 X 是可度量化的（利用第 I 章，§ 9，第 10 号，定理 5 化归为 X 是 σ-紧的情形；然后利用命题 16 的推论）{[}参看 § 5，习题 15{]}。
\end{enumerate}

¶ 24) a) 设 X 是一个完备度量空间，A 是 X 的一个子集，它是一个可数开集族的交。证明 A 上存在一个度量，A 关于它是一个完备度量空间，且它在 A 上既定义由 X 的拓扑诱导的拓扑，又定义一个比 X 的一致结构所诱导的一致结构更细的一致结构{[}参看 § 1，习题 16 b) 与 c){]}。

\begin{enumerate}
\def\labelenumi{\alph{enumi})}
\setcounter{enumi}{1}
\tightlist
\item
  反之，设 X 是一个可度量化空间，A 是 X 的一个子集，使得 A 上存在一个度量 d，它与由 X 的拓扑诱导的拓扑相容，且 A 关于它是一个完备度量空间。证明 A 是 X 中可数个开集的交。{[}对每个整数 n \textgreater{} 0，考虑集合 \(G_{n}\) ，它由具有如下性质的点 \(x \in \overline{A}\) 组成：x 有一个开邻域 U，使 \(A \cap U\) （关于 d）的直径 \(\leqslant r/n\) 。{]}
\end{enumerate}

¶ 25) 设 X 是一个可度量化空间，βX 是它的斯通–切赫紧化（§ 1，习题 7），d 是与 X 的拓扑相容的一个有界度量。对每个 x ∈ X，令 f\_x(y) 表示把 y → d(x, y) 连续地扩张到 βX 上所得的实值函数。

\begin{enumerate}
\def\labelenumi{\alph{enumi})}
\item
  证明 \(f_{x}(z) + f_{y}(z) \geqslant d(x, y)\) 对所有 \(x\) 、\(y\) ∈ \(\mathbf{X}\) 及 \(z\) ∈ \(\beta \mathbf{X}\) 成立。对每个 \(y \in \beta \mathbf{X}\) （异于 \(x \in \beta \mathbf{X}\) ），证明 \(f_{x}(y) > 0\) 。b) 证明若 \(\mathbf{X}\) 关于度量 \(d\) 是完备的，则 \(\mathbf{X}\) 是 \(\beta \mathbf{X}\) 中一个开集序列的交。{[}考虑集合 \(G_{n}\) ，它由所有满足下列条件的 \(y \in \beta \mathbf{X}\) 组成：\(f_{x}(y) < 1 / n\) 对至少一个点 \(x \in \mathbf{X}\) 成立，并利用 a) 证明 \(\mathbf{X} = \bigcap_{n=1}^{\infty} G_{n}\) 。{]}
\item
  反之，证明若 X 是一个开集序列 \(G_{n}\) （含于 \(\beta X\) ）的交，则 X 上存在一个度量 \(d'\) ，它与 X 的拓扑相容，且 X 关于它是完备的。{[}只需考虑 \(X \neq \beta X\) 的情形，注意 \(\beta X - X\) 是紧集组成的一个族 \(F_{n} = \beta X - G_{n}\) 的并；对每个 n，令 \(g_{n}(x) = \inf_{y \in F_{n}} f_{x}(y)\) ，并利用 a) 证明 \(g_{n}(x) > 0\) 对所有 \(x \in X\) 成立，且 \(g_{n}\) 在 X 上连续。证明的结尾按 § 1，习题 16 b) 的方式进行。{]}
\end{enumerate}

